\chapter{ws-Closed Sets and ws-Continuous Maps in Bitopological Spaces}
\graphicspath{{Chapter7/Chapter7Figs/EPS/}{Chapter7/Chapter7Figs/}}

\section{Introduction}\label{sec7.1}

The triple $(\TSP, \tau_1, \tau_2)$ where $\tau_1$ and $\tau_2$ are topologies and P is a set is called a bitopological space. In 1963 Kelly [45] started the investigation of such spaces. He generalized the topological concepts to bitopological setting and published a large number of papers. After Kelly’s work on bitopological spaces, various authors, like Arya and Nour [7], Di Maio and Noiri [24], Fukutake [34], Nagaveni [65], Maki, Sundaram and Balachandran [55], Sheik John [85], Sampath Kumar[82], Patty [77], Arockiarani [3], Gnanambal [40], Reilly [8], Rajamani and Viswanthan [80], and Popa [78] have made stremndous effort on all concepts of topology to bitopological spaces instead of topological spaces.

In chapter II, we have introduced and analysed the concept of ws-closed sets, ws-open sets and ws-closure operator in topological spaces. In chapter III, we have introduced and investigated some properties of ws-continuous maps and ws-irresolute maps in topological spaces.

In section 2 of this chapter, $(\TSi,\TSj)$-ws-closed sets in bitopological space have been introduced and studied. In section 3 of this chapter, we have introduced $(\TSi, \TSj)$-ws-open sets in bitopological space and study some of their related properties. In section 4 of this chapter, we shall use the $(\TSi, \TSj)$-ws-closed subsets of bitopological space $(\TSP, \tau_1, \tau_2)$ to define a new closure operator ``$(\TSi, \TSj)$-ws-cl'', and thus new topology $\tau_{\ws}(\TSi, \TSj)$ on the space and shall examine some of the properties of this new topology.

In section 5 of this chapter, a new kind of maps called $\clrD_{\ws} (\TSi, \TSj)-\sigma_{\TSk}$-continuous maps in bitopological spaces are introduced and investigated. During this process, some of their properties are obtained. Also, we have introduced the concept of ws-bi-continuity, ws-s-bi-continuity and examine some related properties.

Throughout this chapter $(\TSP, \tau_1, \tau_2)$, $(\TSY, \sigma_{1}, \sigma_{2} )$ and $(\TSZ, \eta_{1}, \eta_{2})$ denote nonempty bitopological spaces on which no separation axioms are assumed, unless otherwise mentioned and fixed integers $\TSi, \TSj, \TSk, \TSe, \TSm, \TSn\in \{1, 2\}$.

\section{\boldmath$(\tau_{\TSi}, \tau_{\TSj})$ ws-closed sets and their basic properties}

Here we work on concept of $(\tau_{\TSi}, \tau_{\TSj})$-ws-closed sets which are developed on a bitopological space in analogy with ws-closed sets in topological spaces. From now on, $\tau$-$\cl(\clrD)$ denotes the closure of $\clrD$ relative to a topology $\tau$.

\begin{dfn}\label{defi7.2.1}
Take up $\TSi$, $\TSj \in \{1, 2\}$ be fixed integers. In a bitopological space $(\TSP, \tau_1, \tau_2)$, a subset $\clrD \subset \TSP$ is said to be $(\tau_{\TSi}, \tau_{\TSj})$-ws-closed (briefly, $(\TSi, \TSj)$-ws-closed) set if $\tau_{\TSj}$-$\scl(\clrD) \subset \TSG$ and $G \in \TSw\TSO(\TSP, \tau_{\TSi})$.
\end{dfn}

We denote the family of all $(\TSi, \TSj)$ ws-closed sets in a bitopological space $(\TSP, \tau_1, \tau_2)$ by $\clrD_{\ws}(\tau_{\TSi}, \tau_{\TSj})$ or $\clrD_{\ws}(\TSi, \TSj)$.

\begin{rem}\label{rem7.2.2}
By setting $\tau_1 = \tau_2$ in Definition \ref{defi7.2.1}, an $(\TSi, \TSj)$-ws-closed set reduces to a ws-closed set in $\TSP$.
\end{rem}

Firstly we express that the class of $(\TSi, \TSj)$-ws-closed sets properly lies between the class of $(\TSi, \TSj)$-g\#-closed sets and the class of $(\TSi, \TSj)$-gspr-closed sets.

\begin{thm}\label{thm7.2.3}
If $\TSM$ is $(\TSi, \TSj)$-$\TSg\#$-closed subset of $(\TSP, \tau_1, \tau_2)$, then $\TSM$ is $(\TSi, \TSj)$-$\TSw$-closed.
\end{thm}

\begin{proof}
Take up $\TSM$ be a $(\TSi, \TSj)$-ws-closed subset of $(\TSP, \tau_1, \tau_2)$. Assume $\TSG\in \TSW\TSO(\TSP, \tau_{\TSi})$ be $\ni: \TSA\subset \TSG$. Seeing that $\TSW\TSO(\TSP, \tau_{\TSi}) \subset \TSS\TSO(\TSP, \tau_{\TSi})$, hence we have $\TSG\in\TSS\TSO(\TSP, \tau_{\TSi})$. Then by hypothesis, $\tau_{\TSj}$-$\scl(\TSM) \subset \TSG$. Henceforth $\TSM$ is $(\TSi, \TSj)$-g\#-closed.

We use example \ref{exm7.2.4} to prove the inverse of theorem is untrue.
\end{proof}

\begin{exm}\label{exm7.2.4}
Take up $\TSP= \{1,2,3\}$, $\tau_1 =\{\TSP, \phi, \{1\}, \{2,3\}\}$ and $\tau_2 =\{\TSP, \phi, \{1\},\break \{2\}, \{1,2\}\}$. Then the subsets $\{1\}$ and $\{2\}$ are $(1, 2)$-ws-closed sets, but not $(1, 2)$-g\#-closed sets in the bitopological space $(\TSP, \tau_1, \tau_2 )$.
\end{exm}

\begin{thm}\label{thm7.2.5}
If $\TSA$ is $(\TSi, \TSj)$-$*\TSg\alpha$-closed subset of $(\TSX, \tau_1, \tau_2)$, then $\TSA$ is $(\TSi, \TSj)$-$\TSw$-closed.
\end{thm}

\begin{proof}
Let $\TSA$ be a $(\TSi, \TSj)$-ws-closed subset of $(\TSX, \tau_1, \tau_2)$. Let $\TSG\in\TSW\TSO(\TSX, \tau_{\TSi})$ be such that $\TSA \subset \TSG$. Since $\TSW\TSO(\TSX, \tau_{\TSi}) \subset \TSS\TSO(\TSX, \tau_{\TSi})$, we have $\TSG\in\TSS\TSO(\TSX, \tau_{\TSi})$. Then by hypothesis, $\tau_{\TSj}$-$\scl(\TSA) \subset \TSG$. Henceforth $\TSA$ is $(\TSi, \TSj)$-g\#-closed.

We use example \ref{exm7.2.6} to prove the inverse of theorem is untrue.
\end{proof}

\begin{exm}\label{exm7.2.6}
Let $\TSX= \{1, 2, 3\}$, $\tau_1 =\{\TSX, \phi, \{1\}, \{2,3\}\}$ and $\tau_2 =\{\TSX, \phi, \{1\}, \{2\},\break \{1,2\}\}$. Then the subset $\{2\}$ is $(1, 2)$-ws-closed sets, but not $(1, 2)$-g\#-closed sets in the bitopological space $(\TSX, \tau_1, \tau_2 )$.
\end{exm}

\begin{thm}\label{thm7.2.7}
If $\TSM$ is a $(\TSi, \TSj)$-$\TSg\xi*$-closed subset of $(\TSP, \tau_1, \tau_2)$, then $\TSM$ is $(\TSi, \TSj)$-ws-closed.
\end{thm}

\begin{proof}
Take up $\TSM$ be a $(\TSi, \TSj)$-ws-closed subset of $(\TSP, \tau_1, \tau_2 )$. Assume $\TSG\in \TSW\TSO(\TSP, \tau_{\TSi})$ be $\ni: \TSM \subset
\TSG$. Seeing that $\TSW\TSO(\TSP, \tau_{\TSi}) \subset \TSS\TSO(\TSP, \tau_{\TSi})$, hence we have $\TSG\in\TSS\TSO(\TSP, \tau_{\TSi})$. Then by hypothesis, $\tau_{\TSj}$-$\scl(\TSM) \subset \TSG$. Henceforth $\TSM$ is a $(\TSi, \TSj)$-$\TSg\xi*$-closed.

We use example \ref{exm7.2.8} to prove the inverse of theorem is untrue.
\end{proof}

\begin{exm}\label{exm7.2.8} 
Take up $\TSP= \{1,2,3\}$, $\tau_1 =\{\TSP, \phi, \{1\}, \{2,3\}\}$ and $\tau_2 =\{\TSP, \phi, \{1\},\break \{2\}, \{1, 2\}\}$. Then the subset $\{2\}$ is $(1, 2)$-ws-closed sets, but not $(1, 2)$-$\TSg\xi*$-closed sets in the bitopological space $(\TSP, \tau_1, \tau_2 )$.
\end{exm}

\begin{thm}\label{thm7.2.9}
If $\TSM$ is a $(\TSi, \TSj)$-$\alpha gp$-closed subset of $(\TSP, \tau_1, \tau_2)$, then $\TSM$ is $(\TSi, \TSj)$-ws-closed.
\end{thm}

\begin{proof}
Take up $\TSM$ be a $(\TSi, \TSj)$-ws-closed subset of $(\TSP, \tau_1, \tau_2)$. Assume $\TSG\in\TSW\TSO(\TSP, \tau_{\TSi})$ be $\ni: \TSM \subset \TSG$. Seeing that $\TSW\TSO(\TSP, \tau_{\TSi}) \subset \TSS\TSO(\TSP, \tau_{\TSi})$, hence we have $\TSG\in\TSS\TSO(\TSP, \tau_{\TSi})$. Then by hypothesis, $\tau_{\TSj}$-$
\scl(\TSM) \subset \TSG$. Henceforth $\TSM$ is a $(\TSi, \TSj)$-$\alpha gp$-closed.

We use example \ref{exm7.2.10} to prove the inverse of theorem is untrue.
\end{proof}

\begin{exm}\label{exm7.2.10}
Take up $\TSP= \{1,2,3\}$, $\tau_1 =\{\TSP, \phi, \{1\}, \{2,3\}\}$ and $\tau_2 =\{\TSP, \phi, \{1\},\break \{2\}, \{1, 2\}\}$. Then the subset $\{2\}$ is $(1, 2)$-ws-closed sets, but not $(1, 2)$-$\alpha gp$-closed sets in the bitopological space $(\TSP, \tau_1, \tau_2 )$.
\end{exm}

\begin{thm}\label{thm7.2.11}
If $\TSM$ is a $(\TSi, \TSj)$-$\ddot{g}$-closed subset of $(\TSP, \tau_1, \tau_2)$, then $\TSM$ is $(\TSi, \TSj)$-ws-closed.
\end{thm}

\begin{proof}
Take up $\TSA$ be a $(\TSi, \TSj)$-ws-closed subset of $(\TSP, \tau_1, \tau_2 )$. Assume $\TSG\in\TSW\TSO(\TSP, \tau_{\TSi} )$ be $\ni: \TSM \subset \TSG$. Seeing that $\TSW\TSO(\TSP, \tau_{\TSi} ) \subset \TSS\TSO(\TSP, \tau_{\TSi} )$, hence we have $\TSG\in\TSS\TSO(\TSP, \tau_{\TSi} )$. Then by hypothesis, $\tau_{\TSj}$-$\scl(\clrD) \subset \TSG$. Henceforth $\clrD$ is a $(\TSi, \TSj)$-$\ddot{g}$-closed.

We use example \ref{exm7.2.12} to prove the inverse of theorem is untrue.
\end{proof}

\begin{exm}\label{exm7.2.12}
Take up $\TSP= \{1,2,3\}$, $\tau_1 =\{\TSP, \phi, \{1\}, \{2, 3\}\}$ and $\tau_2 =\{\TSP, \phi, \{1\},\break \{2\}, \{1, 2\}\}$. Then the subset $\{2\}$ is $(1, 2)$-ws-closed sets, but not $(1, 2)$-$\ddot{g}$-closed sets in the bitopological space $(\TSP, \tau_1, \tau_2 )$.
\end{exm}

\begin{thm}\label{thm7.2.13}
If $\clrD$ is a $(\TSi, \TSj)$-rb -closed subset of $(\TSP, \tau_1, \tau_2 )$, then $\clrD$ is $(\TSi, \TSj)$-ws-closed.
\end{thm}

\begin{proof}
Take up $\clrD$ be a $(\TSi, \TSj)$-ws-closed subset of $(\TSP, \tau_1, \tau_2)$. Assume $\TSG\in\TSW\TSO(\TSP, \tau_{\TSi})$ be $\ni: \TSM \subset \TSG$. Seeing that $\TSW\TSO(\TSP, \tau_{\TSi} ) \subset \TSS\TSO(\TSP, \tau_{\TSi})$, hence we have $\TSG\in\TSS\TSO(\TSP, \tau_{\TSi})$. Then by hypothesis, $\tau_{\TSj}$-$\scl(\TSM) \subset \TSG$. Henceforth $\TSM$ is a $(\TSi, \TSj)$-rb-closed.

We use example \ref{exm7.2.14} to prove the inverse of theorem is untrue.
\end{proof}

\begin{exm}\label{exm7.2.14}
Take up $\TSP= \{1,2,3\}$, $\tau_1 =\{\TSP, \phi, \{1\}, \{2, 3\}\}$ and $\tau_2 =\{\TSP, \phi, \{1\},\break \{2\}, \{1, 2\}\}$. Then the subset $\{2\}$ is $(1, 2)$-ws-closed sets, but not $(1, 2)$-rb-closed sets in the bitopological space $(\TSP, \tau_1, \tau_2 )$.
\end{exm}

\begin{thm}\label{thm7.2.15}
If $\TSM$ is a $(\TSi, \TSj)$-ws-closed subset of $(\TSP, \tau_1, \tau_2 )$, then $\TSM$ is $(\TSi, \TSj)$-gspr-closed.
\end{thm}

\begin{proof}
Take up $\TSM$ be a $(\TSi, \TSj)$-ws-closed subset of $(\TSP, \tau_1, \tau_2 )$. Assume $\TSG\in\TSW\TSO(\TSP, \tau_{\TSi} )$ be $\ni: \TSM \subset \TSG$. Since $\TSW\TSO(\TSP, \tau_{\TSi} ) \subset \TSS\TSO(\TSP, \tau_{\TSi} )$, hence we have $\TSG\in\TSS\TSO(\TSP, \tau_{\TSi} )$. Then by hypothesis, $\tau_{\TSj}$-$\Scl(\TSM) \subset \TSG$. Also $\tau_{\TSj}$-$\spcl(\TSM) \subset \tau_{\TSj}$-$\Scl(\TSM)$ which implies $\tau_{\TSj}$-$\spcl(\TSM)\subset \TSG$. Therefore $\TSM$ is $(\TSi, \TSj)$-$\gspr$-closed.

We use example \ref{exm7.2.16} to prove the inverse of theorem is untrue.
\end{proof}

\begin{exm}\label{exm7.2.16}
Take up $\TSP=\{1, 2, 3\}$, $\tau_1 =\{\TSP, \phi, \{1\}, \{ 2,3\}\}$ and $\tau_2 =\{\TSP, \phi, \{1\},\break \{2\},\{1,2\}\}$. Then the subsets $\{1,2\}$ are $(1, 2)$-gspr-closed sets but not $(1, 2)$-ws-closed sets in the bitopological space $(\TSP, \tau_1, \tau_2 )$.
\end{exm}

\begin{thm}\label{thm7.2.17}
If $\TSM$ is a $(\TSi, \TSj)$-ws-closed subset of $(\TSP, \tau_1, \tau_2 )$, then $\TSM$ is $(\TSi, \TSj)$-gsp-closed.
\end{thm}

\begin{proof}
Take up $\TSM$ be a $(\TSi, \TSj)$-ws-closed subset of $(\TSP, \tau_1, \tau_2)$. Let $\TSG\in\TSW\TSO(\TSP, \tau_{\TSi} )$ be $\ni: \TSM \subset \TSG$. Since $\TSW\TSO(\TSP, \tau_{\TSi}) \subset \TSS\TSO(\TSP, \tau_{\TSi})$, we have $\TSG\in\TSS\TSO(\TSP, \tau_{\TSi})$. Then by hypothesis, $\tau_{\TSj}$-$\scl(\TSA) \subset \TSG$. Also $\tau_{\TSj}$-$\spcl(\TSM) \subset \tau_{\TSj}$-$\scl(\TSM)$ which implies $\tau_{\TSj}$-$\spcl(\TSM) \subset \TSG$. Henceforth $\TSM$ is $(\TSi, \TSj)$-gsp-closed.

We use example \ref{exm7.2.18} to prove the inverse of theorem is untrue.
\end{proof}

\begin{exm}\label{exm7.2.18}
Take up $\TSP=\{1,2,3\}$, $\tau_1 =\{\TSP, \phi, \{1\}, \{2, 3\}\}$ and $\tau_2 =\{\TSP, \phi, \{1\},\break \{2\},\{1,2\}\}$. Then the subsets $\{1,2\}$ are $(1, 2)$-gsp-closed sets but not $(1, 2)$-ws-closed sets in the bitopological space $(\TSP, \tau_1, \tau_2 )$.
\end{exm}

\begin{thm}\label{thm7.2.19}
If $\TSM$ is a $(\TSi, \TSj)$-ws-closed subset of $(\TSP, \tau_1, \tau_2 )$, then $\TSM$ is $(\TSi, \TSj)$-rgb-closed.
\end{thm}

\begin{proof}
Take up $\TSM$ be a $(\TSi, \TSj)$-ws-closed subset of $(\TSP, \tau_1, \tau_2 )$. Assume $\TSG\in\TSW\TSO(\TSP, \tau_{\TSi} )$ be $\ni: \TSM \subset\TSG$. Seeing that $\TSW\TSO(\TSP, \tau_{\TSi} ) \subset \TSS\TSO(\TSP, \tau_{\TSi} )$, hence we have $\TSG\in\SO(\TSP, \tau_{\TSi} )$. Then by hypothesis, $\tau_{\TSj}$-$\scl(\TSM) \subset \TSG$. Also $\tau_{\TSj}$-$\bcl(\TSM) \subset \tau_{\TSj}$-$\scl(\TSM)$ which implies $\tau_{\TSj}$-$\bcl(\TSM) \subset \TSG$. Therefore $\TSM$ is $(\TSi, \TSj)$-rgb-closed.

We use example \ref{exm7.2.20} to prove the inverse of theorem is untrue.
\end{proof}

\begin{exm}\label{exm7.2.20}
Take up $\TSP=\{1,2,3\}$, $\tau_1 =\{\TSP, \phi, \{1\}, \{ 2,3\}\}$ and $\tau_2 =\{\TSP, \phi, \{1\},\break \{2\},\{1,2\}\}$. Then the subsets $\{1, 2\}$ are $(1, 2)$-rgb-closed sets but not $(1, 2)$-ws-closed sets in the bitopological space $(\TSP, \tau_1, \tau_2 )$.
\end{exm}

\begin{rem}\label{rem7.2.21}
$\tau_{\TSj}$-rw-closed($\TSR$*-closed)sets, and $(\TSi,\TSj)$-ws-closed sets are independent as seen from the following examples.
\end{rem}

\begin{exm}\label{exm7.2.22}
Take up $\TSP= \{1,2,3\}$, $\tau_1 =\{\TSP, \phi, \{1\}, \{2,3\}\}$ and $\tau_2 =\{\TSP, \phi, \{1\},\break \{2\},\{1, 2\}\}$. Then the subset $\{2\}$ is $(1, 2)$-ws-closed sets, but not $\tau_2$-rw-closed($\TSR$*-closed)sets set in the bitopological space $(\TSP, \tau_1, \tau_2 )$.
\end{exm}

\begin{exm}\label{exm7.2.23}
Take up $\TSP=\{1,2,3\}$, $\tau_1 =\{\TSP, \phi, \{1\}, \{2\}, \{1, 2\}\}$ and $\tau_2 =\{\TSP, \phi, \{1\},\{2\},\{1,2\},\{1, 3\}\}$. Then the subsets $\{1,2\}$ are $\tau_2$-rw-closed($\TSR$*-closed)sets but not $(1, 2)$-ws-closed sets in the bitopological space $(\TSP, \tau_1, \tau_2 )$.
\end{exm}

\begin{rem}\label{rem7.2.24}
From the above discussions and known facts we have the following implications:
\begin{figure}[H]
\centering
\includegraphics[scale=.75]{1.eps}
\end{figure}
\end{rem}

\begin{thm}\label{thm7.2.24}
If $\TSM$, $\TSN\in \clrD_{\ws} (\TSi, \TSj)$, then $\TSM\cap \TSN\in \clrD_{\rw} (\TSi, \TSj)$.
\end{thm}

\begin{proof}
Take up $\TSG\in \TSW\TSO(\TSP, \tau_{\TSi})$ be $\ni: \TSM\cap \TSN \subset \TSG$. Then $\TSM \subset \TSG$ and $\TSN \subset \TSG$. Seeing that $\TSM$, $\TSN\in \clrD_{\ws} (\TSi, \TSj)$, we have $\tau_{\TSj}$-$\scl(\TSM) \subset \TSG$ and $\tau_{\TSj}$-$\scl(\TSN) \subset \TSG$. That is $\tau_{\TSj}$-$\scl(\TSM) \cup \tau_{\TSj}$-$\scl(\TSN) \subset \TSG$. Also $\tau_{\TSj}$-$\scl(\TSM) \cap  \tau_{\TSj}$-$\scl(\TSN)= \tau_{\TSj}$-$\cl(\TSM\cap \TSN)$ and so $\tau_{\TSj}$-$\scl(\TSM\cap \TSN) \subset \TSG$. Therefore $\TSM\cap \TSN\in \clrD_{\ws} (\TSi, \TSj)$.
\end{proof}

\begin{rem}\label{rem7.2.25}
The union of two $(\TSi, \TSj)$-ws-closed sets is generally not a $(\TSi, \TSj)$-ws-closed set as seen from example given below.
\end{rem}

\begin{exm}\label{exm7. 2.26}
Take up $\TSP=\{1,2,3\}$, $\tau_1 =\{\TSP, \phi, \{1\},\{2, 3\}\}$ and $\tau_2 =\{\TSP, \phi, \{1\},\break \{2\}, \{1, 2\}\}$. Then subsets $\{1\}$ and $\{2\}$ are $(1, 2)$-ws-closed sets, but $\{1\}\cup \{2\}=\{1,2\}$ is not a $(1, 2)$-ws-closed set in the bitopological space $(\TSP, \tau_1, \tau_2 )$.
\end{exm}

\begin{rem}\label{rem7.2.27} 
As usual, The family $\clrD_{\ws} (1, 2)$ is unequal to the family $\clrD_{\ws} (2, 1)$ as observed in the next example.
\end{rem}

\begin{exm}\label{exm7.2.28}
Take up $\TSP= \{1,2,3\}$, $\tau_1 =\{\TSP, \phi, \{1\}, \{2\}, \{1, 2\}\}$ and $\tau_2 =\{\TSP, \phi, \{3\}, \{2, 3\}\}$. Then $\clrD_{\ws}(1, 2)=\{\TSP, \phi, \{1\}, \{3\},\{2,3\},\{1,3\}\}$ and $\clrD_{\ws} (2, 1)=\{\TSP, \phi, \{1\}, \{2\}, \{3\},\{2, 3\}, \{1,3\}\}$. Therefore $\clrD_{\ws} (1, 2)\neq \clrD_{\ws} (2, 1)$.
\end{exm}

\begin{thm}\label{thm7.2.29}
If $\tau_1 \subset \tau_2$ and $\SWO(\TSP, \tau_1) \subset \WO(\TSP, \tau_2 )$ in $(\TSP, \tau_1, \tau_2 )$, then $\clrD_{\ws} (\tau_1, \tau_2 ) \supset  \clrD_{\ws} (\tau_2,\tau_1 )$.
\end{thm}

\begin{proof}
Take up $\TSM\in \clrD_{\ws} (\tau_2, \tau_1 )$. That is $\TSM$ is a $(2, 1)$-ws-closed set. To prove that $\TSM\in \clrD_{\ws} (\tau_1,\tau_2 )$. Let $\TSG\in \WO(\TSP, \tau_1 )$ be $\ni: \TSM \subset \TSG$. Since $\WO(\TSP, \tau_1 ) \subset \WO(\TSP, \tau_2 )$, we have $\TSG\in \WO(\TSP, \tau_2 )$. As $\TSM$ is a $(2, 1)$-ws-closed set, we have $\tau_1$-$\scl(\TSM) \subset \TSG$. Since $\tau_1 \subset \tau_2$, we have $\tau_2$-$\scl(\TSM) \subset \tau_1$-$\scl(\TSM)$ and it follows that $\tau_2$-$\scl(\TSM) \subset \TSG$. Hence $\TSM$ is $(1, 2)$-ws-closed. That is $\TSM \in \clrD_{\ws} (\tau_1, \tau_2 )$. Henceforth $\clrD_{\ws} (\tau_1, \tau_2 ) \supset  \clrD_{\ws} (\tau_2, \tau_1 )$.
\end{proof}

\begin{thm}\label{thm7.2.30}
Take up $\TSi$, $\TSj$ be fixed integers of $\{1, 2\}$. For each $\TSP$ of $(\TSP, \tau_1, \tau_2 )$, $\{P\}$ is a $\TSw$-open in $(\TSP, \tau_{\TSi} )$ or $\{P\}^{\text{c}}$ is $(\TSi, \TSj)$-ws-closed.
\end{thm}

\begin{proof}
Pretend $\{P\}$ is not w-open in $(\TSP, \tau_{\TSi} )$. Then $\{P\}^{\text{c}}$ is not w-open in $(\TSP, \tau_{\TSi} )$. Now w-open in $(\TSP, \tau_{\TSi} )$ containing $\{P\}^{\text{c}}$ is $\TSP$ alone. Also $\tau_{\TSj}$-$\scl(\{\TSP\}^{\text{c}}) \subset \TSP$. Hence $\{P\}^{\text{c}}$ is $(\TSi, \TSj)$-ws-closed.
\end{proof}

\begin{coro}\label{coro7.2.31}
If $\TSM$ is $(\TSi, \TSj)$-ws-closed in $(\TSP, \tau_1, \tau_2 )$, then $\TSM$ is $\tau_{\TSj}$-closed iff $\tau_{\TSj}$-$\scl(\TSM)-\TSM$ is a $\tau_{\TSi}$-$\TSw$-open set.
\end{coro}

\begin{proof}
Pretend $\TSM$ is $\tau_{\TSj}$-closed. Thus $\tau_{\TSj}$-$\scl(\TSM) = \TSM$ and so $\tau_{\TSj}$-$\scl(\TSM)-\TSM=\phi$, is $\tau_{\TSi}$-regular semiopen set.

Inversely, Pretend $\tau_{\TSj}$-$\scl(\TSM)-\TSM$ is a $\tau_{\TSi}$-$\TSw$-open. Since $\TSM$ is $(\TSi, \TSj)$-ws-closed, seeing that $\tau_{\TSj}$-$\scl(\TSM)-\TSM\not\in$ non-empty $\tau_{\TSi}$-w-open set. Henceforth $\tau_{\TSj}$-$\scl(\TSM)-\TSM=\phi$. That is $\tau_{\TSj}$-$\scl(\TSM)=\TSM$ and hence $\TSM$ is $\tau_{\TSj}$-closed.
\end{proof}

\begin{thm}\label{thm7.2.32}
In a bitopological space $(\TSP, \tau_1, \tau_2 )$, $\WO(\TSP, \tau_{\TSi} ) \subset \{\TSF\subset\TSP: \TSF^{\text{c}} \in\tau_{\TSj}\}$ iff each subset of $(\TSP, \tau_1, \tau_2 )$ is a $(\TSi, \TSj)$-ws-closed.
\end{thm}

\begin{proof}
Pretend that $\WO(\TSP, \tau_{\TSi} ) \subset \{\TSF\subset \TSP: \TSF^{\text{c}} \in\tau_{\TSj} \}$. Let $\TSM$ be any subset of $\TSP$. Let $\TSG\in \WO(\TSP, \tau_{\TSi} )$ be $\ni: \TSM \subset \TSG$. Thus $\tau_{\TSj}$-$\scl(\TSG)=\TSG$. also $\tau_{\TSj}$-$\scl(\TSM) \subset \tau_{\TSj}$-$\scl(\TSG)=\TSG$. That is $\tau_{\TSj}$-$\scl(\TSM) \subset \TSG$. henceforth $\TSM$ is a $(\TSi, \TSj)$-ws-closed set.

Inversely, Pretend that each subset of $(\TSP, \tau_1, \tau_2 )$ is a $(\TSi, \TSj)$-ws-closed set. Let $\TSG\in \WO(\TSP,\tau_{\TSi} )$. Since $\TSG \subset \TSG$ and $\TSG$ is $(\TSi, \TSj)$-ws-closed, we have $\tau_{\TSj}$-$\scl(\TSG) \subset \TSG$. Thus $\tau_{\TSj}$-$\scl(\TSG) =\TSG$ and so $\TSG$ is $\tau_{\TSj}$-closed. That is $\TSG\in\{\TSF\subset \TSP: \TSF^{\text{c}} \in\tau_{\TSj} \}$. Henceforth $\WO(\TSP, \tau_{\TSi} ) \subset \{\TSF\subset \TSP: \TSF^{\text{c}} \in\tau_{\TSj} \}$. 
\end{proof}

\begin{thm}\label{thm7.2.33}
Pretend $\TSM$ be a $(\TSi, \TSj)$-ws-closed subset of a bitopological space $(\TSP, \tau_1, \tau_2 )$. If $\TSM$ is $\tau_{\TSi}$-w-open, then $\TSM$ is $\tau_{\TSj}$-closed.
\end{thm}

\begin{proof}
Take up $\TSM$ be $\tau_{\TSi}$-w-open. Now $\TSM\subset \TSM$. Then by hypothesis $\tau_{\TSj}$-$\scl(\TSM) \subset \TSM$. Therefore $\tau_{\TSj}$-$\scl(\TSM) = \TSM$. That is $\TSM$ is $\tau_{\TSj}$-closed.
\end{proof}

\begin{thm}\label{thm7.2.34}
If $\TSM$ is a $(\TSi, \TSj)$-ws-closed set and $\tau_{\TSi} \subset \WO(\TSP, \tau_{\TSi} )$, then $\tau_{\TSj}$-$\scl(\{\TSP\}) \cap \TSM \neq \phi$ for each $\TSP \in \tau_{\TSj}$-$\scl(\TSM)$.
\end{thm}

\begin{proof}
Take up $\TSM$ be $(\TSi, \TSj)$-ws-closed and $\tau_{\TSi} \subset \WO(\TSP, \tau_{\TSi} )$. Suppose $\tau_{\TSi}$-$\scl(\{P\}) \cap \TSM = \phi$ for some $\TSP \in \tau_{\TSj}$-$\scl(\TSM)$, then $\TSM \subset (\tau_{\TSi}\text{-scl}(\{P\}))^{\text{c}}$. Now $(\tau_{\TSi}\text{-scl}(\{P\}))^{\text{c}} \in\tau_{\TSi} \subset \WO(\TSP, \tau_{\TSi} )$, by hypothesis. That is $(\tau_{\TSi}\text{-scl}(\{P\}))^{\text{c}}$ is $\tau_{\TSi}$-regular semiopen. Since $\TSM$ is $(\TSi, \TSj)$-ws-closed, Hence we have $\tau_{\TSj}\text{-scl}(\TSM) \subset (\tau_{\TSi}\text{-scl}(\{P\}))^{\text{c}}$. This shows that $\TSP \not\in \tau_{\TSj}\text{-scl}(\TSM)$. This contradicts the assumption.
\end{proof}

\begin{thm}\label{thm7.2.35}
If $\TSM$ is a $(\TSi, \TSj)$-ws-closed set and $\TSM \subset \TSN \subset \tau_{\TSj}\text{-scl}(\TSM)$, then $\TSN$ is $(\TSi, \TSj)$-ws-closed.
\end{thm}

\begin{proof}
Take up $\TSG$ be a $\tau_{\TSi}$-w-open set $\ni: \TSN \subset \TSG$. As $\TSM$ is a $(\TSi, \TSj)$-ws-closed set and $\TSM \subset \TSG$, we have $\tau_{\TSj}\text{-scl}(\TSM) \subset \TSG$. Now $\TSN \subset \tau_{\TSj}\text{-scl}(\TSM)$ which implies, $\tau_{\TSj}\text{-scl}(\TSN) \subset \tau_{\TSj}\text{-scl}\{\tau_{\TSj}\text{-scl}(\TSM)\}= \tau_{\TSj}\text{-scl}(\TSM) \subset \TSG$. Thus $\tau_{\TSj}\text{-scl}(\TSN) \subset \TSG$. Therefore $\TSN$ is a $(\TSi, \TSj)$-ws-closed set.
\end{proof}

\begin{coro}\label{coro7.2.36}
Let $\TSK$ and $\TSL$ be subsets of space $(\TSP, \tau_1, \tau_2 )$. If $\TSL$ is $(\TSi, \TSj)$-ws-open and $\TSK \supset  \tau_{\TSj}\text{-int}(\TSL)$, then $\TSK\cap \TSL$ is $(\TSi, \TSj)$-ws-open.
\end{coro}

\begin{proof}
Take up $\TSL$ be $(\TSi, \TSj)$-ws-open and $\TSK \supset \tau_{\TSj}\text{-int}(\TSL)$.That is $\tau_{\TSj}\text{-int}(\TSL)\subset\TSK$. Then $\tau_{\TSj}\text{-int}(\TSL)\subset \TSK\cap \TSL$. Also $\tau_{\TSj}\text{-int}(\TSL) \subset \TSK\cap \TSL\subset\TSL$ and $\TSL$ is $(\TSi, \TSj)$-ws-open. Seeing that $\TSK\cap \TSL$ is also $(\TSi, \TSj)$-ws-open.
\end{proof}

\begin{thm}\label{thm7.2.37}
Each singleton point in space $(\TSP, \tau_1, \tau_2 )$ is neither $(\TSi, \TSj)$-ws-open nor $\tau_{\TSi}$-w-open.
\end{thm}

\begin{proof}
Take up $(\TSP, \tau_1, \tau_2 )$ be a bitopological space. Assume $\TSx\in\TSP$. To prove $\{\TSx\}$ is either $(\TSi, \TSj)$-ws-open or $\tau_{\TSi}$-w-open. That is to prove $\TSP-\{\TSx\}$, is either $(\TSi, \TSj)$-ws-closed or $\tau_{\TSi}$-w-open, which follows from previous results.
\end{proof}

\section{\boldmath$(\tau_{\TSi}, \tau_{\TSj} )$-ws-open sets and their basic properties}

In this section, we introduce $(\TSi, \TSj)$-ws-open sets in bitopological spaces and examine some of their properties.

\begin{dfn}\label{defi7.3.1}
Let $\TSi$, $\TSj\in\{1, 2\}$ be fixed integers. In a bitopological space $(\TSP, \tau_1, \tau_2 )$, a subset $\TSD \subset \TSP$ is said to be $(\tau_{\TSi}, \tau_{\TSj} )$-ws-open (briefly, $(\TSi, \TSj)$-ws-open) if $\clrD^{\text{c}}$ is $(\TSi, \TSj)$-ws-closed.
\end{dfn}

\begin{thm}\label{thm7.3.2}
In a bitopological space $(\TSP, \tau_1, \tau_2 )$, we have the following
\begin{enumerate}[(1)]
\item Each $(\TSi, \TSj)$-$\TSg\#$-open set is $(\TSi, \TSj)$-ws-open but inverse is untrue.
\item Each $(\TSi, \TSj)$-$\TSg\xi*$-open set is $(\TSi, \TSj)$-ws-open but inverse is untrue.
\item Each $(\TSi, \TSj)$-$\alpha gp$-open set is $(\TSi, \TSj)$-ws-open but inverse is untrue.
\item Each $(\TSi, \TSj)$-$\ddot{g}$-open set is $(\TSi, \TSj)$-rw-open but inverse is untrue.
\item Each $(\TSi, \TSj)$-rb-open set is $(\TSi, \TSj)$-rw-open but inverse is untrue.
\item Each $(\TSi, \TSj)$-ws-open set is $(\TSi, \TSj)$-gspr-open but inverse is untrue.
\item Each $(\TSi, \TSj)$-ws-open set is $(\TSi, \TSj)$-gsp-open but inverse is untrue.
\item Each $(\TSi, \TSj)$-ws-open set is $(\TSi, \TSj)$-rgb-open but inverse is untrue.
\end{enumerate}
\end{thm}

\begin{proof}
Proof is obtained from the Theorems \ref{thm7.2.3}, to \ref{thm7.2.18}.
\end{proof}

\begin{thm}\label{thm7.3.3}
If $\clrD$ and $\TSE$ are $(\TSi, \TSj)$-ws-open sets, then $\TSD\cup \TSE$ is $(\TSi, \TSj)$-ws-open.
\end{thm}

\begin{proof}
The result is obtained from theorem \ref{thm7.2.25}.
\end{proof}

\begin{rem}\label{rem7.3.4}
The intersection of two $(\TSi, \TSj)$-ws-open sets is generally not an $(\TSi, \TSj)$-ws-open set as seen from example \ref{exam7.3.5}.
\end{rem}

\begin{exm}\label{exam7.3.5}
Take $\TSP=\{ 1, 2, 3\}$, $\tau_1 =\{\TSP, \phi, \{1\}, \{2,3\}\}$ and $\tau_2 =\{\TSP, \phi, \{1\}, \{2\},\break \{1, 2\}\}$. Then the subsets $\{1, 3\}$ and $\{1, 3\}$ are $(1, 2)$-ws-open sets, but $\{1, 3\}\cap \{2,3\}=\{3\}$ is not $(1, 2)$-ws-open set in the bitopological space $(\TSP, \tau_1, \tau_2 )$.
\end{exm}

\begin{thm}\label{thm7.3.6}
$\TSA subset \clrD$ of $(\TSP, \tau_1, \tau_2 )$ is $(\TSi, \TSj)$-ws-open iff $\TSF\subset \tau_{\TSj}\text{-int}(\clrD)$, whenever $\TSF$ is $\tau_{\TSi}$-w-open and $\TSF\subset \clrD$.
\end{thm}

\begin{proof}
Pretend that $\TSF\subset \tau_{\TSi}\text{-int}(\clrD)$ whenever $\TSF\subset \clrD$ and $\TSF$ is $\tau_{\TSi}$-w-open. To show that $\clrD$ is $(\TSi, \TSj)$-ws-open. Assume that $\TSG$ be $\tau_{\TSi}$-w-open and $\clrD^{\text{c}} \subset \TSG$. Thun $\TSG^{\text{c}} \subset \clrD$ and $\TSG^{\text{c}}$ is $\tau_{\TSi}$-w-open, by Lemma \ref{lem1.2.5}. By hypothesis, $\TSG^{\text{c}} \subset \tau_{\TSj}\text{-int}(\clrD)$. That is $(\tau_{\TSj}\text{-int}(\clrD))^{\text{c}} \subset \TSG$, since $\tau_{\TSj}\text{-scl}(\clrD^{\text{c}})=(\tau_{\TSj}\text{-int}(\clrD))^{\text{c}}$. Thus $\TSD^{\text{c}}$ is $(\TSi, \TSj)$-ws-closed. That is $\clrD$ is $(\TSi, \TSj)$-ws-open. 

Inversly, Pretend that $\clrD$ is $(\TSi, \TSj)$-ws-open, $\TSF\subset\clrD$ and $\TSF$ is $\tau_{\TSi}$-w-open. Then $\clrD^{\text{c}} \subset \TSF^{\text{c}}$ and $\TSF^{\text{c}}$ is also $\tau_{\TSi}$-w-open by Lemma \ref{lem1.2.5}. Seeing that $\clrD^{\text{c}}$ is $(\TSi, \TSj)$-ws-closed, we have $\tau_{\TSj}\text{-scl}(D^{\text{c}}) \subset \TSF^{\text{c}}$ and so $\TSF\subset \tau_{\TSj}\text{-int}(\clrD)$, since $\tau_{\TSj}\text{-scl}(\clrD^{\text{c}}) = (\tau_{\TSj}\text{-int}(\clrD))^{\text{c}}$.
\end{proof}


\section{\boldmath$(\tau_{\TSi}, \tau_{\TSj} )$-ws-closure in bitopological spaces}\label{sec7.3}

W. Dunham [33] developed the concept of generalized closure operator $\TSC$* and Fukutake [34] derived and examined the concept of pairwise generalized closure operator $(\tau_{\TSi},\tau_{\TSj} )$-$\cl$* in a bitopological spaces. Analogous to that we introduce the pairwise ws-closure operator ``$(\TSi, \TSj)$-ws-$\cl$'' in bitopological spaces and examine some of their related properties.

\begin{dfn}\label{defi7.3.1}
Let $(\TSP, \tau_1, \tau_2 )$ be a bitopological space and $\TSi$, $\TSj\in\{1, 2\}$ be fixed integers. For each subset $\TSE$ of $\TSP$, define $(\tau_{\TSi}, \tau_{\TSj} )$-$\ws$-$\cl(\TSE)= \cap \{\TSG: \TSE \subset \TSG\in\clrD \ws (\TSi, \TSj)\}$ (briefly: $(\TSi, \TSj)$-$\ws$-$\cl(\TSE))$.
\end{dfn}

\begin{thm}\label{thm7.3.2}
If $\TSG$ and $\TSH$ be subsets of $(\TSP, \tau_1, \tau_2 )$. Then
\begin{enumerate}[(1)]
\item $(\TSi, \TSj)$-$\ws$-$\cl(\TSP)=\TSP$ and $(\TSi, \TSj)$-$\ws$-$\cl(\phi)=\phi$.
\item $\TSG \subset (\TSi, \TSj)$-$\ws$-$\cl(\TSG)$.
\item If $\TSH$ is any $(\TSi, \TSj)$-$\ws$-closed set containing $\TSG$, then $(\TSi, \TSj)$-$\ws$-$\cl(\TSG)\subset\TSH$.
\end{enumerate}
\end{thm}

\begin{proof}
Obtained form the Definition \ref{defi7.3.1}.
\end{proof}

\begin{thm}\label{thm7.3.3}
Let $\TSG$ and $\TSH$ be subsets of $(\TSP, \tau_1, \tau_2 )$ and $\TSi$, $\TSj\in\{1, 2\}$ be fixed integers. If $\TSG \subset \TSH$, then $(\TSi, \TSj)$-$\ws$-$\cl(\TSG) \subset (\TSi, \TSj)$-$\ws$-$\cl(\TSH)$.
\end{thm}

\begin{proof}
Take up $\TSG\subset\TSH$. By Definition \ref{defi7.3.1}, $(\TSi, \TSj)$-$\ws$-$\cl(\TSH)=\cap \{\TSF:H\subset\TSF\in\clrD_{\ws}(\TSi, \TSj)\}$ If $H\subset\TSF\in\clrD_{\ws} (\TSi,\TSj)$, since $\TSG\subset\TSH$, $\TSG\subset\TSH\subset\TSF\in\clrD_{\ws} (\TSi, \TSj)$, we have $(\TSi, \TSj)$-$\ws$-$\cl(\TSG) \subset \TSF$. Henceforth $(\TSi, \TSj)$-$\ws$-$\cl(\TSG) \subset \cap \{\TSF:\TSH\subset\TSF\in\clrD_{\ws} (\TSi, \TSj)\}= (\TSi, \TSj)$-$\ws$-$\cl(\TSH)$. That is $(\TSi, \TSj)$-$\ws$-$\scl(\TSG) \subset (\TSi, \TSj)$-$\ws$-$\cl(\TSH)$.
\end{proof}

\begin{thm}\label{thm7.3.4}
Let $\TSG$ be a subset of $(\TSP, \tau_1, \tau_2 )$. If $\tau_1, \subset \tau_2$ and $\WO(\TSP, \tau_1 ) \subset \WO(\TSP, \tau_2 )$, then $(1,2)$-$\ws$-$\cl(\TSG)\subset (2, 1)$-$\ws$-$\cl(\TSG)$.
\end{thm}

\begin{proof}
By definition \ref{defi7.3.1}, $(1,2)$-$\ws$-$\cl(\TSG)=\cap \{\TSF:\TSG\subset F\in\clrD_{\ws} (1, 2)\}$. Seeing that $\tau_1, \subset \tau_2$, by Theorem \label{thm7.2.34} $\clrD_{\ws} (2, 1) \subset \clrD_{\ws} (1, 2)$. Therefore $\cap \{\TSF:\TSG\subset \TSF\in\clrD_{\ws} (1, 2)\}\subset \cap \{\TSF:\TSG\subset \TSF\in\clrD_{\ws} (2,1)\}$. That is $(1, 2)$-$\ws$-$\cl(\TSG)\subset \cap \{\TSF:\TSG\subset \TSF\in\clrD_{\ws} (2, 1)\}= (2, 1)$-$\ws$-$\cl(\TSG)$. Hence $(1, 2)$-$\ws$-$\cl(\TSG)\subset (2, 1)$-$\ws$-$\cl(\TSG)$.
\end{proof}

\begin{thm}\label{thm7.3.5}
Let $\TSG$ be a subset of $(\TSP, \tau_1, \tau_2 )$ and $\TSi$, $\TSj\in\{1, 2\}$ be fixed integers, then $\TSG \subset (\TSi, \TSj)$-$\ws$-$\cl(\TSG) \subset \tau_{\TSj}$-$\scl(\TSG)$.
\end{thm}

\begin{proof}
By Definition \ref{defi7.3.1}, it is obtained that $\TSG \subset (\TSi, \TSj)$-$\ws$-$\cl(\TSG)$. Now to prove that $(\TSi, \TSj)$-$\ws$-$\cl(\TSG) \subset \tau_{\TSj}$-$\scl(\TSG)$. By definition of closure, $\tau_{\TSj}\text{-scl}(\TSG)=\cap \{\TSF\subset\TSP: \TSG\subset\TSF$ and $\TSF$ is $\tau_{\TSj}$-closed$\}$. If $\TSG \subset \TSF$ and $\TSF$ is $\tau_{\TSj}$-closed, then $\TSF$ is $(\TSi, \TSj)$-ws-closed, as each $\tau_{\TSj}$-closed set is is $(\TSi, \TSj)$-ws-closed. Henceforth $(\TSi, \TSj)$-$\ws$-$\cl(\TSG) \subset \cap \{\TSF\subset\TSP: \TSG\subset \TSF$ and $\TSF$ is $\tau_{\TSj}\text{-closed}= \tau_{\TSj}\text{-scl}(\TSG)$. Hence $(\TSi, \TSj)$-$\ws$-$\cl(\TSG)\subset \tau_{\TSj}\text{-scl}(\TSG)$.
\end{proof}

\section{\boldmath$\clrD_{\ws} (\TSi, \TSj)$-$\sigma_{\TSk}$-continuous maps and\hfill\break ws-bi-conti\-nuous maps}

In this section a new class of maps called $\clrD_{\ws} (\TSi, \TSj)$-$\sigma_{\TSk}$-continuous maps in bitopological spaces are developed and examined. During this process, some of their related properties are obtained. Also, we form the concept of ws-bi-continuity and ws-s-bi-continuity in bitopological spaces and examine some of their related properties.

\begin{dfn}\label{defi7.4.1}
A map $\TSh: (\TSP, \tau_1, \tau_2 )\to (\TSQ, \sigma_{1}, \sigma_{2} )$ is called $\clrD_{\ws} (\TSi, \TSj)$-$\sigma_{\TSk}$-continuous if the inverse image of every $\sigma_{\TSk}$-closed set is an $(\TSi, \TSj)$-ws-closed set in $(\TSP, \tau_1, \tau_2 )$.
\end{dfn}

\begin{rem}\label{rem7.4.2}
If $\tau_1 =\tau_2 =\tau$ and $\sigma_{1} =\sigma_{2} =\sigma$  in Definition \ref{defi7.4.1}, then the $\clrD_{\ws} (\TSi, \TSj)$-$\sigma_{\TSk}$-continuity of maps coincides with ws-continuity of maps in topological spaces.
\end{rem}

\begin{thm}\label{thm7.4.3}
If a map $\TSh :(\TSP, \tau_1, \tau_2 )\to (\TSQ, \sigma_{1}, \sigma_{2} )$ is $\tau_{\TSj}-\sigma_{\TSk}$-continuous, then it is a $\clrD_{\ws} (\TSi, \TSj)-\sigma_{\TSk}$-continuous.
\end{thm}

\begin{proof}
Take $\TSL$ be a $\sigma_{\TSk}$-closed set. Since $\TSh$ is $\tau_{\TSj}-\sigma_{\TSk}$-continuous, $\TSh^{-1} (\TSL)$ is $\tau_{\TSj}$-closed. By Theorem \ref{thm7.2.7}, $\TSh^{-1} (\TSL)$ is $(\TSi, \TSj)$-ws-closed in $(\TSP, \tau_1, \tau_2 )$. Therefore $\TSh$ is $\clrD_{\ws}(\TSi, \TSj)-\sigma_{\TSk}$-continuous.

We use example \ref{exam7.4.4} to prove the inverse of theorem is untrue.
\end{proof}

\begin{exm}\label{exma7.4.4}
Let $\TSP=\{1, 2, 3, 4\}$, $\tau_1 =\{\TSP, \phi, \{1\},\{2\},\{1,2\},\{1,2,3\}$ and $\tau_2 =\{\TSP, \phi, \{1\},\{1, 2\}\{1,2, 3\}\}$, $\TSQ=\{1, 2, 3\}$, $\sigma_{1} =\{\TSQ, \phi, \{1\},\{2\},\{1,2\}\}$ and $\sigma_{2} =\{\TSQ, \phi, \{1\},\{2,3\}\}$. Define a map $\TSh: (\TSP,\tau_1, \tau_2 )\to (\TSQ, \sigma_{1}, \sigma_{2} )$ by $\TSh(1)=4$, $\TSh(2)=2$, $\TSh(3)=2$, $\TSh(4)=3$. Then $\TSh$ is $\clrD_{\rw}(1, 2)$-$\sigma_{2}$-continuous but it is not $\tau_1-\sigma_{2}$-continuous, Seeing that for the $\sigma_{2}$-closed set $\{2,3\}$, $\TSh^{-1}(\{2,3\})=\{2, 3\}$ is not $\tau_1$-closed.
\end{exm}

\begin{thm}\label{thm7.4.5}
Map $\TSh: (\TSP, \tau_1, \tau_2 )\to (\TSQ, \sigma_{1}, \sigma_{2} )$ is $(\TSi, \TSj)$-$\sigma_{\TSk}$-$\TSg\#$-continuous, then it is $\clrD_{\ws} (\TSi, \TSj)$-$\sigma_{\TSk}$-continuous.
\end{thm}

\begin{proof}
Let $\TSL$ be a $\sigma_{\TSk}$-closed set. Since $\TSh$ is $(\TSi, \TSj)$-$\sigma_{\TSk}$-$\TSg\#$-continuous, $\TSh^{-1}(\TSL)$ is $(\TSi, \TSj)$-$\TSg\#$-closed. By Theorem \ref{thm7.2.3}, $\TSh^{-1}(\TSL)$ is $(\TSi, \TSj)$-ws-closed in $(\TSP, \tau_1, \tau_2 )$. Therefore $\TSh$ is $\clrD_{\ws} (\TSi, \TSj)$-$\sigma_{\TSk}$-continuous.

We use example \ref{exam7.4.6} to prove the inverse of theorem is untrue.
\end{proof}

\begin{exm}\label{exam7.4.6}
Let $\TSP=\{1, 2, 3\}$, $\tau_1 =\{\TSP, \phi, \{1\},\{2,3\}\}$ and $\tau_2 =\{\TSP, \phi, \{1\},\{2\},\break \{1, 2\}\}$ $\TSQ=\{1, 2,3\}$, $\sigma_{1} =\{\TSQ,\phi,\{1\},\{2\},\{1,2\},\{1,3\}\}$ and $\sigma_{2} =\{\TSQ,\phi,\{1\}\},(1,2)-\clrD_{\ws}=\{\TSP,\phi,\{1\},\{2\},\{3\}\{2,3\},\{1,3\}\}$, and $(1,2)$-$\TSg\#=\{\TSP, \phi, \{3\},\{2,3\},\{1,2\}\}$, Define a map $\TSh: (\TSP, \tau_1,\tau_2 )\to (\TSQ, \sigma_{1}, \sigma_{2} )$ by $\TSh(1)=1$, $\TSh(2)=2$, $\TSh(3)=3$. Then $\TSh$ is $\clrD_{\ws}(1, 2)-\sigma_{1}$-continuous but it is not $(\TSi,\TSj)$-$\sigma_{\TSk}$-$\TSg\#$-continuous. Seeing that for the $\sigma_{1}$-closed set $\{2\}$, $\TSh^{-1}(\{2\})=\{2\}$ is not $(1, 2)$-$\TSg\#$-closed set.
\end{exm}

\begin{thm}\label{thm7.4.7}
Map $\TSh: (\TSP, \tau_1, \tau_2 )\to (\TSQ, \sigma_{1}, \sigma_{2} )$ is $(\TSi, \TSj)$-$\sigma_{\TSk}$-$\TSg\xi*$-continuous, then it is $\clrD_{\ws} (\TSi, \TSj)$-$\sigma_{\TSk}$-continuous.
\end{thm}

\begin{proof}
Assume $\TSL$ be a $\sigma_{\TSk}$-closed set. Seeing that $\TSh$ is $(\TSi, \TSj)$-$\sigma_{\TSk}$-$\TSg\xi*$-continuous, $\TSh^{-1}(\TSL)$ is $(\TSi, \TSj)$-$\TSg\xi*$-closed. By Theorem \ref{thm7.2.7}, $\TSh^{-1}(\TSL)$ is $(\TSi, \TSj)$-ws-closed in $(\TSP, \tau_1, \tau_2 )$. Hnceforth $\TSh$ is $\clrD_{\ws} (\TSi, \TSj)$-$\sigma_{\TSk}$-continuous.

We use example \ref{exam7.4.8} to prove the inverse of theorem is untrue.
\end{proof}

\begin{exm}\label{exam7.4.8}
Let $\TSP=\{1,2,3\}$, $\tau_1 =\{\TSP, \phi, \{1\},\{2,3\}$ and $\tau_2 =\{\TSP, \phi, \{1\},\{2\},\break \{1,2\}\}$ $\TSQ=\{1,2,3\}$, $\sigma_{1} =\{\TSQ, \phi,\{1\},\{2\},\{1,2\},\{1,2\}\}$ and $\sigma_{2} =\{\TSQ,\phi,\{2\}\}$, $(1,2)- \clrD_{\ws} =\{\TSP, \phi, \{1\},\{2\},\{3\}\{2,3\},\{1,3\}\}$, and $(1,2)$-$\TSg\xi*=\{ \TSP, \phi, \{1\},\{3\},\break \{2,3\},\{1,2\}\}$, Define a map $\TSh: (\TSP, \tau_1, \tau_2 )\to (\TSQ, \sigma_{1}, \sigma_{2} )$ by $\TSh(1)=1$, $\TSh(2)=2$, $\TSh(3)=3$. Then $\TSh$ is $\clrD_{\ws} (1, 2)$-$\sigma_{1}$-continuous but it is not $(\TSi, \TSj)$-$\sigma_{\TSk}$-$\TSg\xi*$-continuous. Seeing that for the $\sigma_{1}$-closed set $\{2\}$, $h^{-1}(\{2\})=\{2\}$ is not $(1, 2)-\TSg\xi*$-closed set.
\end{exm}

\begin{thm}\label{thm7.4.9}
Map $\TSh :(\TSP, \tau_1, \tau_2 )\to (\TSQ, \sigma_{1}, \sigma_{2} )$ is $(\TSi, \TSj)$-$\sigma_{\TSk}$-$\alpha gp$-continuous, then it is $\clrD_{\ws} (\TSi, \TSj)$-$\sigma_{\TSk}$-continuous.
\end{thm}

\begin{proof}
Take up $\TSL$ be a $\sigma_{\TSk}$-closed set. See $\TSh$ is $(\TSi, \TSj)$-$\sigma_{\TSk}$-$\alpha gp$-continuous, $\TSh^{-1} (\TSL)$ is $(\TSi, \TSj)$-$\alpha gp$-closed. By Theorem \ref{thm7.2.9}, $\TSh^{-1} (\TSL)$ is $(\TSi, \TSj)$-ws-closed in $(\TSP, \tau_1, \tau_2 )$. Hencefoth $\TSh$ is $\clrD_{\ws} (\TSi, \TSj)$-$\sigma_{\TSk}$-continuous.

We use example \ref{exam7.4.10} to prove the inverse of theorem is untrue.
\end{proof}

\begin{exm}\label{exam7.4.10}
Let $\TSP=\{1,2,3\}$, $\tau_1 =\{\TSP, \phi, \{1\},\{2,3\}$ and $\tau_2 =\{\TSP, \phi, \{1\},\{2\},\break \{1,2\}\}$ $\TSQ=\{1,2,3\}$, $\sigma_{1} =\{\TSQ,\phi,\{1\},\{2\},\{1,2\},\{1,3\}\}$ and $\sigma_{2} =\{\TSQ,\phi, \{2\}\}$, $(1,2)$-$\clrD_{\ws} =\{\TSP, \phi, \{1\},\{2\},\{3\}\{2,3\},\{1,3\}\}$, and $(1,2)$-$\alpha gp =\{\TSP, \phi, \{1\},\{2,3\}\}$, Define a map $\TSh: (\TSP, \tau_1, \tau_2 )\to (\TSQ, \sigma_{1}, \sigma_{2} )$ by $\TSh(1)=1$, $\TSh(2)=2$, $\TSh(3)=3$. Then h is $\clrD_{\ws} (1, 2)$-$\sigma_{1}$-continuous but it is not $(\TSi, \TSj)$-$\sigma_{\TSk}$-$\alpha gp$-continuous. Seeing that for the $\sigma_{1}$-closed set $\{2\}$, $\TSh^{-1} (\{2\})=\{2\}$ is not $(1, 2)$-$\alpha gp$ closed set.
\end{exm}

\begin{thm}\label{thm7.4.11}
Map $\TSh:(\TSP, \tau_1, \tau_2 )\to (\TSQ, \sigma_{1}, \sigma_{2} )$ is $(\TSi, \TSj)$-$\sigma_{\TSk}$-$\ddot{g}$-continuous, then it is $\clrD_{\ws} (\TSi, \TSj)$-$\sigma_{\TSk}$-continuous.
\end{thm}

\begin{proof}
Take $\TSL$ be a $\sigma_{\TSk}$-closed set. Seeing that $\TSh$ is $(\TSi, \TSj)$-$\sigma_{\TSk}$-$\ddot{g}$-continuous, $\TSh^{-1} (\TSL)$ is $(\TSi, \TSj)$-$\ddot{g}$-closed. By Theorem \ref{thm7.2.11}, $\TSh^{-1} (\TSL)$ is $(\TSi, \TSj)$-ws-closed in $(\TSP, \tau_1, \tau_2 )$. Henceforth $\TSh$ is $\clrD_{\ws} (\TSi, \TSj)$-$\sigma_{\TSk}$-continuous.

We use example \ref{exam7.4.12} to prove the inverse of theorem is untrue.
\end{proof}

\begin{exm}\label{exam7.4.12}
Let $\TSP=\{1,2,3\}$, $\tau_1 =\{\TSP, \phi, \{1\},\{2,3\}$ and $\tau_2 =\{\TSP, \phi, \{1\},\{2\},\break \{1,2\}\}$ $\TSQ=\{1,2,3\}$, $\sigma_{1} =\{\TSQ,\phi,\{1\},\{2\},\{1,2\},\{1,2\}\}$ and $\sigma_{2} =\{\TSQ,\phi, \{2\}\}$, $(1,2)$-$\clrD_{\ws} =\{\TSP, \phi, \{1\},\{2\},\{3\}\{2,3\},\{1,3\}\}$, and $(1,2)$-$\ddot{g} =\{\TSP, \phi, \{1\},\{2,3\}\}$, Define a map $\TSh: (\TSP, \tau_1, \tau_2 )\to (\TSQ, \sigma_{1}, \sigma_{2} )$ by $\TSh(1)=1$, $\TSh(2)=2$, $\TSh(3)=3$. Then $\TSh$ is $\clrD_{\ws} (1, 2)$-$\sigma_{1}$-continuous but it is not $(\TSi, \TSj)$-$\sigma_{\TSk}$-$\ddot{g}$-continuous. Seeing that for the $\sigma_{1}$-closed set $\{2\}$, $\TSh^{-1} (\{2\})=\{2\}$ which is not $(1, 2)$-$\ddot{g}$ closed set.
\end{exm}

\begin{thm}\label{thm7.4.13}
Map $\TSh:(\TSP, \tau_1, \tau_2 )\to (\TSQ, \sigma_{1}, \sigma_{2} )$ is $(\TSi, \TSj)$-$\sigma_{\TSk}$-rb-continuous, then it is $\clrD_{\ws} (\TSi, \TSj)$-$\sigma_{\TSk}$-continuous.
\end{thm}

\begin{proof}
Let $\TSL$ be a $\sigma_{\TSk}$-closed set. Since $\TSh$ is $(\TSi, \TSj)$-$\sigma_{\TSk}$-rb-continuous, $\TSh^{-1} (\TSL)$ is $(\TSi, \TSj)$-rb-closed. By Theorem \ref{thm7.2.13}, $\TSh^{-1} (\TSL)$ is $(\TSi, \TSj)$-ws-closed in $(\TSP, \tau_1, \tau_2 )$. Henceforth $\TSh$ is $\clrD_{\ws} (\TSi, \TSj)$-$\sigma_{\TSk}$-continuous.

We use example \ref{exam7.4.14} to prove the inverse of theorem is untrue.
\end{proof}

\begin{exm}\label{exam7.4.14}
Let $\TSP=\{1,2,3\}$, $\tau_1 =\{\TSP, \phi, \{1\},\{2,3\}$ and $\tau_2 =\{\TSP, \phi, \{1\},\{2\},\break \{1,2\}\}$ $\TSQ=\{1,2,3\}$, $\sigma_{1} =\{\TSQ, \phi,\{1\},\{2\},\{1,2\},\{1,2\}\}$ and $\sigma_{2} =\{\TSQ,\phi,\{2\}\}$, $(1,2)$-$\clrD_{\ws} =\{\TSP, \phi, \{1\},\{2\},\{3\}\{2,3\},\{1,2\}\}$, and $(1,2)$-$\rb=\{\TSP, \phi, \{1\},\{2,3\}\}$, Define a map $\TSh: (\TSP, \tau_1, \tau_2 )\to (\TSQ, \sigma_{1}, \sigma_{2} )$ by $\TSh(1)=1$, $\TSh(2)=2$, $\TSh(3)=3$. Then $\TSh$ is $\clrD_{\ws} (1, 2)$-$\sigma_{1}$-continuous but it is not $(\TSi, \TSj)$-$\sigma_{\TSk}$-rb-continuous. Seeing that for the $\sigma_{1}$-closed set $\{2\}$, $\TSh^{-1} (\{2\})=\{2\}$ which is not $(1, 2)$-rb closed set.
\end{exm}

\begin{thm}\label{thm7.4.15}
Map $\TSh:(\TSP, \tau_1, \tau_2 )\to (\TSQ, \sigma_{1}, \sigma_{2} )$ is $(\TSi, \TSj)$-$\sigma_{\TSk}$-$*\TSg\alpha$-continuous, then it is $\clrD_{\ws}(\TSi, \TSj)$-$\sigma_{\TSk}$-continuous.
\end{thm}

\begin{proof}
Let $\TSL$ be a $\sigma_{\TSk}$-closed set. Since $\TSh$ is $(\TSi, \TSj)$-$\sigma_{\TSk}$-$*\TSg\alpha$-continuous, $\TSh^{-1} (\TSL)$ is $(\TSi, \TSj)$-rb-closed. By Theorem \ref{thm7.2.5}, $\TSh^{-1}(\TSL)$ is $(\TSi, \TSj)$-ws-closed in $(\TSP, \tau_1, \tau_2 )$. Henceforth $\TSh$ is $\clrD_{\ws} (\TSi, \TSj)$-$\sigma_{\TSk}$-continuous.

We use example \ref{exam7.4.16} to prove the inverse of theorem is untrue.
\end{proof}

\begin{exm}\label{exam7.4.16}
Let $\TSP=\{1,2,3\}$, $\tau_1 =\{\TSP, \phi, \{1\},\{2,3\}$ and $\tau_2 =\{\TSP, \phi, \{1\},\{2\},\break \{1,2\}\}$ $\TSQ=\{1,2,3\}$, $\sigma_{1} =\{\TSQ,\phi,\{1\},\{2\},\{1,2\},\{1,2\}\}$ and $\sigma_{2} =\{\TSQ,\phi, \{2\}\}$, $(1,2)$-$\clrD_{\ws} =\{\TSP, \phi, \{1\},\{2\},\{3\}\{2,3\},\{1,3\}\}$, and $(1,2)$-$*\TSg\alpha =\{\TSP, \phi, \{1\},\{2,3\}\}$, Define a map $\TSh: (\TSP, \tau_1, \tau_2 )\to (\TSQ, \sigma_{1}, \sigma_{2} )$ by $\TSh(1)=1$, $\TSh(2)=2$, $\TSh(3)=3$. Then $\TSh$ is $\clrD_{\ws} (1, 2)$-$\sigma_{1}$-continuous but it is not $(\TSi, \TSj)$-$\sigma_{\TSk}$-$*\TSg\alpha$-continuous, Seeing that for the $\sigma_{1}$-closed set $\{2\}$, $\TSh^{-1} (\{2\})=\{2\}$ is not $(1, 2)$-$*\TSg\alpha$-closed set.
\end{exm}

\begin{thm}\label{thm7.4.17}
Map $\TSh:(\TSP, \tau_1, \tau_2 )\to (\TSQ, \sigma_{1}, \sigma_{2} )$ is $\clrD_{\ws} (\TSi, \TSj)$-$\sigma_{\TSk}$-continuous, then it is $(\TSi, \TSj)$-$\sigma_{\TSk}$-$\gspr$-continuous.
\end{thm}

\begin{proof}
Take $\TSL$ be a $\sigma_{\TSk}$-closed set. Since $\TSh$ is $\clrD_{\ws}(\TSi, \TSj)$-$\sigma_{\TSk}$-continuous, $\TSh^{-1}(\TSL)$ is $(\TSi, \TSj)$-ws-closed. By Theorem \ref{thm7.2.15}, $\TSh^{-1} (\TSL)$ is $(\TSi, \TSj)$-$\gspr$-closed in $(\TSP, \tau_1, \tau_2 )$. Henceforth $\TSh$ is $(\TSi, \TSj)$-$\sigma_{\TSk}$-$\gspr$-continuous.

We use example \ref{exam7.4.18} to prove the inverse of theorem is untrue.
\end{proof}

\begin{exm}\label{exam7.4.18}
Let $\TSP=\{1,2,3\}$, $\tau_1 =\{\TSP, \phi, \{1\}, \{2,3\}\}$ $\tau_2 =\{\TSP, \phi, \{1\},\{2\},\{1,2\}\}$ and $\TSQ=\{1,2,3\}$, $\sigma_{1} =\{\TSQ, \phi, \{3\}\}$ and $\sigma_{2} =\{\TSQ, \phi\}$. $(1, 2)$-$\clrD_{\ws} =\{\TSP, \phi, \{1\},\{2\},\{3\},\break \{2,3\},\{1,3\}\}$, and $(1,2)$-$\gspr =\{\TSP(\TSP)\}$, Define a map $\TSh:(\TSP, \tau_1, \tau_2 )\to (\TSQ, \sigma_{1}, \sigma_{2} )$ by $\TSh(1)=1$, $\TSh(2)=2$ and $\TSh(3)=3$, Then $\TSh$ is $(1, 2)$-$\sigma_{1}$-$\gspr$-continuous but it is not $\clrD_{\ws} (1, 2)$-$\sigma_{1}$-continuous. Seeing that for the $\sigma_{1}$-closed set $\{1,2\}$, $\TSh^{-1} (\{1,2\})=\{1,2\}$ is not a $(1, 2)$-ws-closed set.
\end{exm}

\begin{thm}\label{thm7.4.19}
Map $\TSh:(\TSP, \tau_1, \tau_2 )\to (\TSQ, \sigma_{1}, \sigma_{2} )$ is $\clrD_{\ws} (\TSi, \TSj)$-$\sigma_{\TSk}$-continuous, then it is $(\TSi, \TSj)$-$\sigma_{\TSk}$-$\gsp$-continuous.
\end{thm}

\begin{proof}
Let $\TSL$ be a $\sigma_{\TSk}$-closed set. Since $\TSh$ is $\clrD_{\ws} (\TSi, \TSj)$-$\sigma_{\TSk}$-continuous, $\TSh^{-1} (\TSL)$ is $(\TSi, \TSj)$-ws-closed. By Theorem \ref{thm7.2.17}, $\TSh^{-1} (\TSL)$ is $(\TSi, \TSj)$-$\gsp$-closed in $(\TSP, \tau_1, \tau_2 )$. Henceforth $\TSh$ is $(\TSi, \TSj)$-$\sigma_{\TSk}$-$\gsp$ continuous.

We use example \ref{exam7.4.20} to prove the inverse of theorem is untrue.
\end{proof}

\begin{exm}\label{exam7.4.20}
Let $\TSP=\{1,2,3\}$, $\tau_1 =\{\TSP, \phi, \{1\}, \{2,3\}\}$ $\tau_2 =\{\TSP, \phi, \{1\},\{2\},\{1,2\}\}$ and $\TSQ=\{1,2,3\}$, $\sigma_{1} =\{\TSQ, \phi, \{3\}\}$ and $\sigma_{2} =\{\TSQ, \phi\}$. $(1,2)$-$\clrD_{\ws} =\{\TSP, \phi, \{1\},\{2\},\{3\},\break\{2,3\},\{1,3\}\}$, and $(1,2)$-$\gsp =\{\TSP(\TSP)\}$, Define a map $\TSh:(\TSP, \tau_1, \tau_2 )\to (\TSQ, \sigma_{1}, \sigma_{2} )$ by $\TSh(1)=1$, $\TSh(2)=2$ and $\TSh(3)=3$, Then $\TSh$ is $(1, 2)$-$\sigma_{1}$-$\gsp$-continuous but it is not $\clrD_{\ws} (1, 2)$-$\sigma_{1}$-continuous. for the $\sigma_{1}$-closed set $\{1,2\}$, $\TSh^{-1} (\{1,2\})=\{1,2\}$ is not a $(1, 2)$-ws-closed set.
\end{exm}

\begin{thm}\label{thm7.4.21}
Map $\TSh:(\TSP, \tau_1, \tau_2 )\to (\TSQ, \sigma_{1}, \sigma_{2} )$ is $\clrD_{\ws} (\TSi, \TSj)$-$\sigma_{\TSk}$-continuous, then it is $(\TSi, \TSj)$-$\sigma_{\TSk}$-$\rgb$ continuous.
\end{thm}

\begin{proof}
Let $\TSL$ be a $\sigma_{\TSk}$-closed set. Since $\TSh$ is $\clrD_{\ws} (\TSi, \TSj)$-$\sigma_{\TSk}$-continuous, $\TSh^{-1} (\TSL)$ is $(\TSi, \TSj)$-ws-closed. By Theorem \ref{thm7.2.19}, $\TSh^{-1} (\TSL)$ is $(\TSi, \TSj)$-$\rgb$-closed in $(\TSP, \tau_1, \tau_2 )$. Henceforth $\TSh$ is $(\TSi, \TSj)$-$\sigma_{\TSk}$-rgb continuous.

We use example \ref{exam7.4.22} to prove the inverse of theorem is untrue.
\end{proof}

\begin{exm}\label{exam7.4.22}
Let $\TSP=\{1,2,3\}$, $\tau_1 =\{\TSP, \phi, \{1\}, \{2,3\}\}$ $\tau_2 =\{\TSP, \phi, \{1\},\{2\},\break \{1,2\}\}$ and $\TSQ=\{1,2,3\}$, $\sigma_{1} =\{\TSQ, \phi, \{3\}\}$ and $\sigma_{2} =\{\TSQ, \phi\}$. $(1,2)$-$\clrD_{\ws} =\{\TSP, \phi, \{1\},\break \{2\},\{3\}\{2,3\},\{1,2\}\}$, and $(1,2)$-$\rgb =\{\TSP(\TSP)\}$, Define a map $\TSh:(\TSP, \tau_1, \tau_2 )\to (\TSQ, \sigma_{1}, \sigma_{2} )$ by $\TSh(1)=1$, $\TSh(2)=2$ and $\TSh(3)=3$, Then $\TSh$ is $(1, 2)$-$\sigma_{1}$-rgb-continuous but it is not $\clrD_{\ws} (1, 2)$-$\sigma_{1}$-continuous. Seeing that for the $\sigma_{1}$-closed set $\{1,2\}$, $\TSh^{-1} (\{1,2\})=\{1,2\}$ which is not a $(1, 2)$-ws-closed set.
\end{exm}

\begin{rem}\label{rem7.4.23}
$\clrD_{\ws} (\TSi, \TSj)$-$\sigma_{\TSk}$-continuous maps and $\clrD_{\rw} (\TSi, \TSj)$-$\sigma_{\TSk}$-continuous$((\TSi, \TSj)$-$\sigma_{\TSk}$-{\break}$\TSR$*continuous) maps are independent.
\end{rem}

\begin{exm}\label{exam7.4.24}
Let $\TSP=\{1,2,3\}$, $\tau_1 =\{\TSP, \phi, \{1\}, \{2,3\}\}$ and $\tau_2 =\{\TSP, \phi, \{1\},\{2\},\break \{1,2\}\}$ and $\TSQ=\{1,2,3\}$, $\sigma_{1} =\{\TSQ, \phi,\{1\},\{2\},\{1,2\},\{1,2\}\}$ and $\sigma_{2} =\{\TSQ, \phi\}$. Define a map $\TSh:(\TSP, \tau_1, \tau_2 )\to (\TSQ, \sigma_{1},\sigma_{2} )$ by $\TSh(1)=1$,$\TSh(2)=2$ and $\TSh(3)=3$. Then this function $\TSh$ is $\clrD_{\ws} (1, 2)$-$\sigma_{2}$-continuous but it is not $\TSD_{\rw}(1, 2)$-$\sigma_{1}$-continuous$((\TSi, \TSj)$-$\sigma_{\TSk}$-$\TSR*$continuous), Seeing that for the $\sigma_{1}$-closed set $\{2\}$, $\TSh^{-1} (\{2\})=\{2\}$ is not $(1, 2)$-rw-closed($\TSR*$-closed) in $(\TSP, \tau_1, \tau_2 )$. 
\end{exm}

\begin{exm}\label{exam7.4.25}
Let $\TSP=\{1,2,3\}$, $\tau_1 =\{\TSP, \phi, \{1\}, \{2\}, \{1,2\}\}$ and $\tau_2 =\{\TSP, \phi, \{1\},\break \{2\},\{1,2\},\{1,3\}\}$ and $\TSQ=\{1,2,3\}$, $\sigma_{1} =\{\TSQ, \phi, \{3\}\}$. and $\sigma_{2} =\{\TSQ, \phi, \}$. Define a map $\TSh:(\TSP, \tau_1, \tau_2 )\to (\TSQ, \sigma_{1}, \sigma_{2} )$ by $\TSh(1)=1$, $\TSh(2)=2$ and $\TSh(3)=3$. Then this function $\TSh$ is $\clrD_{\rw} (1, 2)$-$\sigma_{1}$-continuous($(\TSi, \TSj)$-$\sigma_{1}$-$\TSR*$continuous) but it is not $\clrD_{\ws} (1, 2)$-$\sigma_{1}$-continuous. Seeing that for the $\sigma$-closed set $\{1,2\}$, $\TSh^{-1}$ ($\{1,2\})=\{1,2\}$ is not $(1, 2)$-ws-closed in $(\TSP, \tau_1, \tau_2 )$.
\end{exm}

\begin{rem}\label{rem7.4.26}
From the above discussions and known results we have the following implications.
\begin{figure}[H]
\centering
\includegraphics[scale=.75]{2.eps}
\end{figure}
\end{rem}

\begin{thm}\label{thm7.4.28}
The continuous following statements are one and the same.
\begin{enumerate}[(1)]
\item A map $\TSh:(\TSP, \tau_1, \tau_2 )\to (\TSQ, \sigma_{1}, \sigma_{2} )$ is $\clrD_{\ws} (\TSi, \TSj)$-$\sigma_{\TSk}$-continuous.

\item The inverse image of every $\sigma_{\TSk}$-open set in $\TSQ$ is $(\TSi, \TSj)$-ws-open in $\TSP$.
\end{enumerate}
\end{thm}

\begin{proof}
(1) $\Rightarrow$  (2) Assume $\TSL$ be a $\sigma_{\TSk}$-open set in $\TSQ$. Then $\TSL^{\text{C}}$ is $\sigma_{\TSk}$-closed set in $\TSQ$. Seeing that $\TSh$ is $\clrD_{\ws}(\TSi, \TSj)$-$\sigma_{\TSk}$-continuous, $\TSh^{-1} (\TSL^{\text{C}})$ is $(\TSi, \TSj)$-ws-closed in $\TSP$. That is $\TSh^{-1} (\TSL^{\text{C}})=(\TSh^{-1} (\TSL))^{\text{C}}$ and so $\TSh^{-1} (\TSL)$ is $(\TSi, \TSj)$-ws-open in $(\TSP, \tau_1, \tau_2 )$.

(2) $\Rightarrow$  (1) Let $\TSM$ be a $\sigma_{\TSk}$-closed set in $\TSQ$. Then $\TSL^{\text{C}}$ is $\sigma_{\TSk}$-open set in $\TSQ$. By hypothesis, $\TSh^{-1} (\TSM^{\text{C}})$ is $(\TSi, \TSj)$-ws-open in $\TSP$. That is $\TSh^{-1} (\TSM^{\text{C}}) = (\TSh^{-1} (\TSM))^{\text{C}}$ and so $\TSh^{-1} (\TSM)$ is $(\TSi, \TSj)$-ws-closed in $(\TSP, \tau_1, \tau_2 )$. Henceforth $\TSh$ is $\clrD_{\ws}(\TSi, \TSj)$-$\sigma_{\TSk}$-continuous.
\end{proof}

\begin{thm}\label{thm7.4.29}
Map $\TSh:(\TSP, \tau_1, \tau_2 )\to (\TSQ, \sigma_{1}, \sigma_{2} )$ is $\clrD_{\ws} (\TSi, \TSj)$-$\sigma_{\TSk}$-continuous and $\TSg: (\TSQ, \sigma_{1}, \sigma_{2} )\to (\TSR, \eta_{1}, \eta_{2})$ is $\sigma_{\TSk}$-$\eta_{\TSn}$-continuous, then $\goh$ is $\clrD_{\ws} (\TSi, \TSj)$-$\eta_{\TSn}$-continuous.
\end{thm}

\begin{proof}
Let $\TSK$ be $\eta_{\TSn}$-closed set in $(\TSR, \eta_{1}, \eta_{2} )$. Seeing that $\TSg$ is $\sigma_{\TSk}$-$\eta_{\TSn}$-continuous, $\TSg^{-1}(\TSK)$ is a $\sigma_{\TSk}$-closed set in $(\TSQ, \sigma_{1}, \sigma_{2} )$. Seeing that $\TSh$ is $\clrD_{\rw}(\TSi, \TSj)$-$\sigma_{\TSk}$-continuous, $\TSh^{-1} (\TSg^{-1}(\TSK))= (\goh)^{-1} (\TSK)$ is a $(\TSi,\TSj)$-ws-closed set in $(\TSP, \tau_1, \tau_2 )$ and hence $\goh$ is $\clrD_{\ws} (\TSi, \TSj)$-$\eta_{\TSn}$-continuous.
\end{proof}

\begin{dfn}\label{defi7.4.30}
\begin{enumerate}[(i)]
\item A map $\TSh:(\TSP, \tau_1, \tau_2 )\to (\TSQ, \sigma_{1}, \sigma_{2} )$ is called ws-bi-continuous if $\TSh$ is $\clrD_{\ws} (1,2)$-$\sigma_{2}$-continuous and is $\clrD_{\ws} (2, 1)$-$\sigma_{1}$-continuous.

\item A map $\TSh:(\TSP, \tau_1, \tau_2 )\to (\TSQ, \sigma_{1}, \sigma_{2} )$ is called $\rw$-strongly-$\bi$-continuous (briefly $\rw$-$\TSs$-$\bi$-continuous) if $\TSh$ is $\rw$-$\bi$-continuous, $\clrD_{\rw} (2, 1)$-$\sigma_{2}$-continuous and\break $\clrD_{\rw}(1, 2)$-$\sigma_{1}$-continuous.
\end{enumerate}
\end{dfn}

\begin{thm}\label{thm7.4.31}
Let $\TSh: (\TSP, \tau_1, \tau_2 )\to (\TSQ, \sigma_{1}, \sigma_{2} )$ be a map.
\begin{enumerate}[(1)]
\item If $\TSh$ is bi-continuous then $\TSh$ is $\ws$-$\bi$-continuous.
\item If $\TSh$ is $\TSs$-$\bi$-continuous then $\TSh$ is $\rw$-$\TSs$-$\bi$-continuous.
\end{enumerate}
\end{thm}

\begin{proof}
\begin{enumerate}[(i)]
\item Take $\TSh:(\TSP, \tau_1, \tau_2 )\to (\TSQ, \sigma_{1}, \sigma_{2} )$ be a $\bi$-continuous map. Then $\TSh$ is $\tau_1 -\sigma_{1}$-continuous and $\tau_2$-$\sigma_{2}$-continuous and so by Theorem \ref{thm7.4.3}, $\TSh$ is $\clrD_{\ws} (1, 2)$-$\sigma_{2}$-continuous and $\clrD_{\ws} (2, 1)$-$\sigma_{1}$-continuous. Thus $\TSh$ is $\ws$-$\bi$-continuous.

\item Similar to (i), using Theorem \ref{thm7.4.3}.
\end{enumerate}

We use example \ref{exam7.4.32} to prove the inverse of theorem is untrue.
\end{proof}

\begin{exm}\label{exam7.4.32}
Let $\TSP=\{1,2,3\}$, $\tau_1 =\{\TSP, \phi, \{1\},\{2,3\}\}$ and $\tau_2 =\{\TSP, \phi,\{1\},\{2\}\break \{1,2\}\}$ and $\TSQ=\{1,2,3\}$, $\sigma_{1} =\{\TSQ, \phi, \{1\}\}$ and $\sigma_{2} =\{\TSQ, \phi, \{1\},\{2\},\{1,2\},\{1,3\}\}$. Define a map $\TSh: (\TSP, \tau_1, \tau_2 )\to (\TSQ, \sigma_{1}, \sigma_{2} )$ by $\TSh(1)=1$, $\TSh(2)=2$, $\TSh(3)=3$. Then $\TSh$ is $\ws$-$\bi$-continuous but not $\bi$-continuous.
\end{exm}

\begin{exm}\label{exam7.4.33}
Let $\TSP=\{1,2,3\}$, $\tau_1 =\{\TSP, \phi, \{1\},\{2,3\}\}$ and $\tau_2 =\{\TSP, \phi,\{1\},\{2\}\break \{1,2\}\}$ and $\TSQ=\{1,2,3\}$, $\sigma_{1} =\{\TSQ, \phi, \{2\}\}$ and $\sigma_{2} =\{\TSQ, \phi, \{1\}\}$. Define a map $\TSh: (\TSP, \tau_1, \tau_2)\to (\TSQ, \sigma_{1}, \sigma_{2} )$ by $\TSh(1)=1$, $\TSh(2)=2$, $\TSh(3)=3$. Then $\TSh$ is $\ws$-$\TSs$-$\bi$-continuous but not $\TSs$-$\bi$-continuous.
\end{exm}
